% !TeX spellcheck = pt_BR
\documentclass[a4paper,12pt]{letter}  % tipo do documento
%\usepackage[brazil]{babel}           % define a linguagem padrão
\usepackage{graphicx}                % permite a inserção de figuras
\usepackage[T1]{fontenc}
% \usepackage[latin1]{inputenc}        % para usar caracteres acentuados diretamente
\usepackage[utf8]{inputenc}
\usepackage[hmargin=1.5cm,vmargin=1cm]{geometry}
\usepackage{listings} %para colocar codigos
\usepackage{amsmath}
\usepackage{amssymb}
\usepackage{multicol}
\setlength{\columnsep}{1.5cm}
\setlength{\columnseprule}{0.2pt}

% \usepackage{splitidx}
% \lstset{numbers=left,
% stepnumber=1,
% firstnumber=1,
% numberstyle=\tiny,
% extendedchars=true,
% breaklines=true,
% frame=tb,
% basicstyle=\footnotesize,
% % stringstyle=\ttfamily,
% % showstringspaces=false
% }




% \newcommand{\gabarito}[1]{\textbf{\begin{small} \\ #1 \\ \end{small}}} %mostrar gabarito
\newcommand{\gabarito}[1]{\textbf{\begin{small} \end{small}}} %nao mostrar gabarito


\begin{document}
\thispagestyle{empty}
% \pagestyle{empty}

UNIVERSIDADE FEDERAL DO CEARÁ\\
Departamento de Estatística e Matemática Aplicada\\
Disciplina de Elementos de Análise Combinatória cc0279.\\
Professor: Albert Einstein F. Muritiba.\\
Temas: Análise Combinatória\\[0.2cm]
\textbf{Nome}: \makebox[380pt]{\hrulefill} \vspace{2pt}
%Matr\'icula:\makebox[60pt]{\hrulefill} \vspace{2pt}
\begin{center}
	2º AP
	% -- \textit{2ª Chamada}
	% 2\textordfeminine Avaliação - Segunda Chamada
	% npc 2 %+ nef
	% NEF
\end{center}
%\begin{multicols}{2}
Resolva as questões a seguir, apresentando detalhadamente todos os cálculos, raciocínios e justificativas.
\begin{enumerate}
    \item Defina e diferencie conceitualmente os tipos de agrupamento: permutações (simples, com repetição, circulares e caóticas) e combinações (simples e completas).
	\item Explique em que tipo de problemas de contagem aplicam-se os Lemas de Kaplansky, diferenciando as situações que exigem o uso de cada um dos dois lemas.
	\item Em um restaurante de \textit{fast-food}, os pratos são montados a partir de 16 opções de ingredientes sólidos, 4 de molhos e 4 de massas. De quantos modos distintos pode-se montar um prato contendo 8 ingredientes sólidos, 2 molhos e 1 massa, sabendo que ingredientes e molhos podem ser repetidos?
	
    \item Uma comissão formada por 2 homens e 2 mulheres deve ser escolhida em um grupo de 10 homens e 8 mulheres. De quantos modos diferentes essa comissão pode ser formada se um determinado homem não aceita participar da comissão se nela estiver presente uma determinada mulher?
    
    \item Determine a quantidade de números inteiros de 1 a 1000, inclusive, que são divisíveis por pelo menos dois dos números 2, 3 e 21.


           
\end{enumerate}

% \begin{scriptsize}
	\textbf{Fórmulas:}\\
	$P_n = n!$\\
	$P^{a,b, ...}_n = \frac{n!}{a!b! ...}$\\
	$PC_n = (n-1)!$\\
	$C^p_n = \frac{n!}{p!(n-p)!}$\\
	$CR^{p}_{n} = C^{p}_{n+p-1}$\\
	Kaplansky: $f(n,p) = C^p_{n-p+1}$ , $g(n,p) = \frac{n}{n-p}C^p_{n-p}$ \\
	Permutação caótica: $D_n = n! \left[ \frac{1}{0!} - \frac{1}{1!} + \frac{1}{2!} - \frac{1}{3!} + \dots + (-1)^n \frac{1}{n!}\right] $\\
	Princípio da Inclusão-Exclusão:\\
	$\sharp \bigcup_{i=1}^n A_i = \sum_{i=1}^n (-1)^{i-1} S_i$\\
	$S_0 = \sharp\Omega $ ,
	$S_1 = \sum_{i=1}^n \sharp A_i $ ,
	$S_2 = \sum_{1\leq i < j \leq n}^n \sharp(A_i \cap A_j) $ ,
	$S_3 = \sum_{1\leq i < j < k \leq n}^n \sharp(A_i \cap A_j \cap A_k) $ \dots \\
	Exatamente p: $\sum_{k=0}^{n-p} (-1)^k C^{k}_{k+p} S_{p+k}$\\
	Pelo menos p: $\sum_{k=0}^{n-p} (-1)^k C^{k}_{k+p-1} S_{p+k}$\\
% \end{scriptsize}

\end{document}
